\newpage
\section{Top-Level-Entity für DPRAM und CPU Integration \\ (Aufgabe 9)}

Dieses Kapitel beschreibt die Integration des WS2812-LED-Controllers in das 
bestehende NEORV32-SoPC. Ziel ist es, eine Dual-Port-RAM-basierte 
Peripherie zu schaffen, die über den Wishbone-Bus von der CPU beschrieben 
und gleichzeitig von der Hardware-FSM zur Ansteuerung der 8×8-RGB-Matrix 
ausgelesen werden kann. Hierzu werden das RAMtoWS2812-Modul um Generics 
erweitert, ein DPRAM als Pixelpuffer instanziiert und eine übergeordnete 
Wishbone-Slave-Entity entworfen. Die Anbindung an das Prozessorsystem 
erfolgt durch Erweiterung des Busmultiplexers im Top-Level-Design. 
Abschließend wird die Funktion durch ein C-Testprogramm verifiziert, das 
die Matrix über Memory-Mapped-I/O mit verschiedenen Farbmustern beschreibt.

\subsection{Anpassung von RAMtoWS2812 auf Generic-Parameter}

Das bestehende RAMtoWS2812-Modul wurde von festen Konstanten auf flexible
 Generics umgestellt. Die Parameter ADDR\textunderscore WIDTH, DATA\textunderscore WIDTH, UNUSED\textunderscore BITS 
 und MAX\textunderscore ADDR ermöglichen eine wiederverwendbare Architektur, die sich 
 an verschiedene Speicherbreiten und Protokollanforderungen anpassen 
 lässt. Das Shiftregister arbeitet nun mit reduzierter Bitbreite, wobei 
 die oberen, ungenutzten Bits des Datenvektors ignoriert werden. Dies ist 
 notwendig, da der Wishbone-Bus mit 32 Bit arbeitet, die WS2812-Logik 
 jedoch nur 24 Bit Farbinformation pro Pixel verarbeitet. Die 
Generic-Schnittstelle erlaubt spätere Änderungen ohne Eingriff in den Modulkern.

\inputminted[linenos,frame=single,fontsize=\footnotesize]{vhdl}{A9_Abgabe/iceduino_setup/osflow/boards/iceduino/RAMtoWS2812.vhd}

\inputminted[linenos,frame=single,fontsize=\footnotesize]{vhdl}{A9_Abgabe/iceduino_setup/osflow/boards/iceduino/addresscounter.vhd}

\inputminted[linenos,frame=single,fontsize=\footnotesize]{vhdl}{A9_Abgabe/iceduino_setup/osflow/boards/iceduino/shiftregister.vhd}


\subsection{Dual-Port-RAM}

Das DPRAM stellt die zentrale Speicherbrücke zwischen CPU und 
WS2812-Controller dar. Zwei unabhängige Ports ermöglichen gleichzeitiges 
Schreiben durch die CPU und Lesen durch die LED-Ansteuerlogik. Die 
Schreibseite empfängt 32-Bit-Worte vom Wishbone-Bus, während die 
Leseite 32-Bit-Blöcke an das Shiftregister weitergibt. Die getrennten 
Takt- und Adressleitungen beider Ports verhindern Zugriffskonflikte und 
gewährleisten deterministisches Timing für die LED-Ausgabe. Die 
Speichergröße von 64 Einträgen deckt exakt die 64 Pixel der 8×8-Matrix ab.

\inputminted[linenos,frame=single,fontsize=\footnotesize]{vhdl}{A9_Abgabe/iceduino_setup/osflow/boards/iceduino/ram.vhd}


\subsection{Wishbone-Bridge neorv\textunderscore ws2812}

Die neorv\textunderscore ws2812-Entity kapselt die gesamte WS2812-Peripherie und stellt 
eine standardisierte Wishbone-Slave-Schnittstelle bereit. Sie instanziiert 
das DPRAM sowie den RAMtoWS2812-Controller und verbindet deren interne Signale. 
Die Adressdekodierung aktiviert das Modul bei Zugriffen auf den konfigurierten 
Basisadressbereich. Der Ausgang ws2812\textunderscore out führt den seriellen Datenstrom zur 
LED-Matrix, während start\textunderscore send die Ausgabe der gespeicherten Pixeldaten auslöst. 
Die Struktur folgt dem etablierten Muster der bestehenden Iceduino-Peripheriemodule.

\inputminted[linenos,frame=single,fontsize=\footnotesize]{vhdl}{A9_Abgabe/iceduino_setup/osflow/boards/iceduino/neorv_ws2812.vhd}


\subsection{Integration in neorv32\textunderscore iceduino\textunderscore top}

Das Top-Level-Entity wurde um eine weitere Wishbone-Response-Instanz erweitert. 
Der Adressbereich F0000100 wurde gewählt, da die oberen 24 Bit (F00001) eine einfache 
Maskierung im Busmultiplexer erlauben – die unteren 8 Bit bleiben für Registeroffsets 
innerhalb des Moduls frei. Der Multiplexer wählt bei Treffer auf diesen Bereich die 
Response-Signale der neorv\textunderscore ws2812-Instanz aus und leitet sie an den Prozessor weiter. 
Der WS2812-Ausgang wurde auf io\textunderscore d(2) gelegt, da dieser Pin über den Arduino-Header 
zugänglich ist und später kompatibel zum Multi-Iceduino-Baseboard bleibt. Die 
Taster-Eingabe für start\textunderscore send nutzt btn(0).

\inputminted[linenos,frame=single,fontsize=\footnotesize, firstline=163, lastline=248, breaklines]{vhdl}{A9_Abgabe/iceduino_setup/osflow/boards/iceduino/neorv32_iceduino_top.vhd}

\newpage

\inputminted[linenos,frame=single,fontsize=\footnotesize, firstline=425, lastline=449]{vhdl}{A9_Abgabe/iceduino_setup/osflow/boards/iceduino/neorv32_iceduino_top.vhd}

\subsection{C-Software mit Hardware-Abstraktion}

Die Softwareschicht definiert über den WS2812\textunderscore BASE-Zeiger einen typisierten Zugriff 
auf das DPRAM. Die Headerdatei deklariert die set\textunderscore pixel-Funktion, die 24-Bit-RGB-Werte 
in 32-Bit-Worte packt und an die Speicheradressen der Matrix schreibt. Das Hauptprogramm 
implementiert einen Zustandsautomaten auf Tastendruck, der die gesamte Matrix sequenziell 
in Rot, Grün oder Blau beschreibt. Die Verzögerungsschleifen zwischen den Schreibzugriffen 
dienen der visuellen Beobachtung; in der Praxis würde die Hardware-FSM die Ausgabe autonom 
übernehmen. Die volatile-Qualifizierung der Pointer verhindert Compiler-Optimierungen bei 
Memory-Mapped-I/O-Zugriffen.


\inputminted[linenos,frame=single,fontsize=\footnotesize, firstline=83, lastline=86]{c}{A9_Abgabe/iceduino_setup/sw/lib/include/neorv32_iceduino.h}
\vspace{-4mm}
\begin{center}
\footnotesize\textit{neorv32\textunderscore iceduino.h}
\end{center}
\vspace{4mm}

\newpage

\inputminted[linenos,frame=single,fontsize=\footnotesize]{c}{A9_Abgabe/iceduino_setup/sw/lib/include/ws2812.h}
\vspace{-4mm}
\begin{center}
\footnotesize\textit{ws2812.h}
\end{center}
\vspace{4mm}

\inputminted[linenos,frame=single,fontsize=\footnotesize]{c}{A9_Abgabe/iceduino_setup/sw/lib/source/ws2812.c}
\vspace{-4mm}
\begin{center}
\footnotesize\textit{ws2812.c}
\end{center}
\vspace{4mm}

\inputminted[linenos,frame=single,fontsize=\footnotesize]{c}{A9_Abgabe/iceduino_setup/sw/programs/iceduino_led/main.c}
\vspace{-4mm}
\begin{center}
\footnotesize\textit{main.c}
\end{center}


